%! TeX root = ../thesis.tex
\chapter{Results}

To characterize the implementation of the hybrid representation,
we can create a single-determinant state $\ket{\psi_0}$
and do a power iteration to generate new states
containing multiple determinants
\begin{equation}
    \label{eq:power-iteration}
    \ket{\psi_n}
    = \frac{\ham \ket{\psi_{n-1}}}{\|\ham \ket{\psi_{n-1}}\|}
    = \frac{\ham^n \ket{\psi_0}}{\|\ham^n \ket{\psi_0}\|}.
\end{equation}
This will be relevant when calculating the ground state using the
Lanczos algorithm~\cite{Lanczos1950} which uses a Krylov subspace
\begin{equation}
    \mathcal{K}_n (\ham, \ket{\psi})
    = \spn\{\ket{\psi}, \ham\ket{\psi}, \ham^2\ket{\psi}, \ldots,
    \ham^n\ket{\psi}\}
\end{equation}
for a given operator $\ham$ and state $\ket{\psi}$.
We create a Hamiltonian for the geometry shown in~\cref{subfig:geometry-tri}
using the parameters listed in
\cref{tab:parameters}.
This Hamiltonian is converted to the natural orbital representation
and the starting state is half-filled
as shown in~\cref{subfig:geometry-nat}
\begin{equation}
    \ket{\psi_0} = \fdag_{\up} \cdag_{b\dn}
    \prod_{v_k,\sigma} \cdag_{v_k\sigma} \ket{0}.
\end{equation}
$\cdag_b$ is the creation operator for the mirror bath site.

Performance benchmarks ran on the Vienna Scientific Cluster VSC-5
using the package \textsc{BenchmarkTools}~\cite{BenchmarkTools.jl-2016}.
The minimum time was used as the runtime estimator
following the authors' suggestion.
The system contains AMD EPYC\texttrademark~7713\footnote{
    \url{https://www.amd.com/en/products/cpu/amd-epyc-7713}
}
CPUs with 64 cores and at least \qty{256}{\giga\byte} of memory.
As the implemented code is currently single-threaded
and memory requirements depend on the model complexity,
it can also run on consumer hardware.

\begin{table}[ht]
    \caption{Hamiltonian parameters used for benchmarks.}
    \label{tab:parameters}
    \centering
    \begin{tabular}{l c S}
        % \toprule
        \doubletoprule
        Parameter             & Symbol       & {Value} \\
        \midrule
        Coulomb interaction   & $U$          & 4.0     \\
        Impurity energy level & $\epsilon$   & -2.0    \\
        Bath energy levels    & $\epsilon_k$ & 0.0     \\
        Hybridization         & $V$          & -1.0    \\
        Bath hopping          & $-t_{ij}$    & -1.0    \\
        \doublebottomrule
    \end{tabular}
\end{table}

\section{Numerical precision}

In order to compare results between the bit and hybrid representation
(BR and HR respectively),
we can write functions to convert between them.
We calculate up to 20 iterations, varying the amount of bath sites,
and number of sites inside the determinant of the HR.
After a couple of iterations we can find a mismatch of SDs,
which are associated with amplitudes
$|\alpha_i| \leq 10^{-20}$ which are much smaller than
machine epsilon $\epsilon = 2^{-52} \approx 10^{-16}$,
and generated when two amplitudes add up destructively.
After removing all SDs with weight $|\alpha_i| < \epsilon$,
each amplitude obtained by the two methods can be compared against each other
and verified to fit within a relative tolerance of \num{e-10}.
For the whole wave function the difference is roughly
\begin{equation}
    \|\ket{\psi^\mathrm{BR}} - \ket{\psi^\mathrm{HR}}\|
    \approx O(10^{-15}).
\end{equation}

\section{Hamiltonian}

\begin{figure}[ht]
    \centering
    %! TeX root = ../thesis.tex
\begin{tikzpicture}
    \begin{groupplot}[
            group style={
                    group size=2 by 1,
                    horizontal sep=15mm,
                    ylabels at=edge left,
                },
            width=0.45\textwidth,
            xlabel={$n_b$ (bath sites)},
            ylabel={$t$ (\unit{\second})},
            xtick={100,200,...,300},
            minor xtick={50,150,...,250},
            xminorticks=true,
            legend style={
                    at={(0.05, 0.95)},
                    anchor=north west,
                },
            legend cell align={left},
        ]

        \nextgroupplot[
            xmode=log,
            ymode=log,
            xticklabel=\pgfmathparse{exp(\tick)}\pgfmathprintnumber{\pgfmathresult},
            xticklabel style={
                    /pgf/number format/fixed,
                    /pgf/number format/precision=0,
                },
        ]
        \addplot+[only marks]
        table [x=bath_sites, y=WF] {./data/hamiltonian_creation.csv};
        \addplot+[only marks]
        table [x=bath_sites, y=CIWF2] {./data/hamiltonian_creation.csv};
        \addplot+[only marks]
        table [x=bath_sites, y=CIWF1] {./data/hamiltonian_creation.csv};
        \addplot
        [
            domain=50:300,
            samples=10,
            dashed,
        ]
        {6e-9 * x^3}
        coordinate [pos=0.9] (cubic);
        \node[inner sep=6mm, left] at (cubic.west) {$O(n_b^3)$};
        \addplot
        [
            domain=50:300,
            samples=10,
            dashed,
        ]
        {2.5e-7 * x^2}
        coordinate [pos=0.9] (quadratic);
        \node[inner sep=4mm, below] at (quadratic.south) {$O(n_b^2)$};
        \legend{
            BR,
            HR-2,
            HR-1,
        };

        \nextgroupplot
        \pgfplotsset{cycle list shift=2}
        \addplot+
        [
            only marks,
            forget plot,
        ]
        table [x=bath_sites, y=CIWF1] {./data/hamiltonian_creation.csv};
        \addplot
        [
            domain=50:300,
            samples=10,
            dashed,
        ]
        {1.59e-7 * x + 5.6e-5}
        coordinate [pos=0.9] (linear);
        \node[inner sep=4mm, left] at (linear.west) {$O(n_b)$};
    \end{groupplot}
\end{tikzpicture}

    \caption{
        Time taken to create the Hamiltonian
        in the natural orbital representation
        using different methods.
        The hybrid representation (HR) uses four sites (including impurity)
        in the bit component,
        single (HR-1) and single plus double (HR-2) excitations in $\ham_v$.
        Right panel shows HR-1 using linear axes.
    }
    \label{fig:hamiltonian-creation}
\end{figure}

For the following tests we put four sites (including impurity) into the
bit component and all other sites into the vector component.
\Cref{fig:hamiltonian-creation} shows the time taken to create the Hamiltonian
for both representations.
The old method shows a cubic scaling while
the new method scales by the maximum allowed excitation.

\section{Power iteration}

\begin{figure}[htb!]
    \centering
    %! TeX root = ../thesis.tex
\begin{tikzpicture}
    \begin{groupplot}[
            group style={
                    group name=plots,
                    group size=2 by 2,
                    horizontal sep=15mm,
                    vertical sep=2mm,
                    x descriptions at=edge bottom,
                    ylabels at=edge left,
                },
            width=0.45\textwidth,
            ymode=log,
            xlabel={$n$ (iteration)},
            xmin=-10,
            xmax=110,
            every axis plot post/.append style={only marks},
        ]

        \nextgroupplot[
            ylabel={$t$ (\unit{\second})},
            legend to name=grouplegend,
            legend columns=2,
            transpose legend=true,
            legend style={
                    /tikz/every even column/.append style={column sep=5mm}
                },
        ]
        % WF
        \addplot+ table [x=i, y=t_wf_1_301] {./data/h_psi_1_100.csv};
        \addlegendentry[color=black]{BR, $n_b=301$}
        \addplot+ table [x=i, y=t_wf_1_101] {./data/h_psi_1_100.csv};
        \addlegendentry[color=black]{BR, $n_b=101$}
        % CIWF
        \addplot+ table [x=i, y=t_ciwf_1_301] {./data/h_psi_1_100.csv};
        \addlegendentry[color=black]{HR, $n_b=301$}
        \addplot+ table [x=i, y=t_ciwf_1_101] {./data/h_psi_1_100.csv};
        \addlegendentry[color=black]{HR, $n_b=101$}

        \nextgroupplot
        % WF
        \addplot+ table [x=i, y=t_wf_2_301] {./data/h_psi_2_100.csv};
        \addplot+ table [x=i, y=t_wf_2_101] {./data/h_psi_2_100.csv};
        % CIWF
        \addplot+ table [x=i, y=t_ciwf_2_301] {./data/h_psi_2_100.csv};
        \addplot+ table [x=i, y=t_ciwf_2_101] {./data/h_psi_2_100.csv};

        \nextgroupplot[
            ylabel={$d$ (determinants)},
        ]
        % WF
        \addplot+ table [x=i, y=l_wf_1_301] {./data/h_psi_1_100.csv};
        \addplot+ table [x=i, y=l_wf_1_101] {./data/h_psi_1_100.csv};
        % CIWF
        \addplot+ table [x=i, y=l_ciwf_1_301] {./data/h_psi_1_100.csv};
        \addplot+ table [x=i, y=l_ciwf_1_101] {./data/h_psi_1_100.csv};

        \nextgroupplot
        % WF
        \addplot+ table [x=i, y=l_wf_2_301] {./data/h_psi_2_100.csv};
        \addplot+ table [x=i, y=l_wf_2_101] {./data/h_psi_2_100.csv};
        % CIWF
        \addplot+ table [x=i, y=l_ciwf_2_301] {./data/h_psi_2_100.csv};
        \addplot+ table [x=i, y=l_ciwf_2_101] {./data/h_psi_2_100.csv};

    \end{groupplot}

    \node at (plots c1r1.north east)
    [
        inner sep=0pt,
        anchor=south,
        yshift=2mm,
    ]
    {\ref{grouplegend}};
\end{tikzpicture}

    \caption{
        Power iteration $n$ of~\cref{eq:power-iteration}
        using \numlist{101; 301} bath sites $n_b$.
        Top panels show time to calculate one step $\ham \ket{\psi_n}$
        excluding normalization.
        Bottom panels show number of determinants of $\ket{\psi_n}$,
        which overlap for the hybrid representation.
        Left: single excitation.
        Right: single plus double excitation.
    }
    \label{fig:wave-function-iteration}
\end{figure}

Since normalization is faster than applying the Hamiltonian,
we only time the numerator of~\cref{eq:power-iteration} in benchmarks.
The top panels of~\cref{fig:wave-function-iteration} show the time taken
to calculate one iteration $\ham\ket{\psi_n}$.
We can see that the time for the HR plateaus already at iteration ten.
This flattening can be understood
if we look at the number of determinants at each step,
shown in the bottom panels of~\cref{fig:wave-function-iteration}.
Since most of the information is stored inside vectors and not the
bit component (containing only four sites),
the number of possible determinants is very low
and independent of the number of bath sites.
For a half-filled system we get
36 determinants in a state of no excitation,
96 in a single excitation state,
and 88 in a double excitation state.
Once these determinants are created,
which happens before the tenth iteration,
there is no additional cost,
and therefore constant time.

The BR with 101 bath sites
restricted to singles
plateaus between iteration \numlist{50; 60}.
This happens because we have two chains with 50 bath sites each.
The excitation starts at the beginning of the chain and hops to the next
site at each iteration.
Once we reach the end,
all possible states are generated.

To get the improvement of the new method over the old one,
we can create the ratio of the measured times
which is shown in~\cref{fig:performance-improvement-iteration}.
We can see that the new method restricted to singles
shows an improvement for all measured parameters.
In the case of 101 bath sites the improvement plateaus
between iteration \numlist{50; 60} which we explained before.

For double excitation we see a performance regression for
the first 15 iterations
which correspond to states with less than around \num{e4} SDs.
For states with more determinants
we gain a performance boost with the HR.

\begin{figure}
    \centering
    %! TeX root = ../thesis.tex
\begin{tikzpicture}
    \begin{groupplot}[
            group style={
                    group size=2 by 1,
                    horizontal sep=15mm,
                    ylabels at=edge left,
                },
            width=0.45\textwidth,
            ymode=log,
            xlabel={$n$ (iteration)},
            ylabel={Improvement},
            xmin=-10,
            xmax=110,
        ]

        \nextgroupplot[
            legend style={
                    at={(0.95, 0.05)},
                    anchor=south east
                },
            legend cell align={left},
        ]
        \addplot+[only marks]
        table [x=i, y=1_301] {./data/wave_function_improvement.csv};
        \addlegendentry{$n_b=301$}
        \addplot+[only marks]
        table [x=i, y=1_101] {./data/wave_function_improvement.csv};
        \addlegendentry{$n_b=101$}

        \nextgroupplot
        \addplot+[only marks]
        table [x=i, y=2_301] {./data/wave_function_improvement.csv};
        \addplot+[only marks]
        table [x=i, y=2_101] {./data/wave_function_improvement.csv};
        % unity
        \draw[dashed] (axis cs:-10,1) -- (axis cs:110,1);
    \end{groupplot}
\end{tikzpicture}

    \caption{
        Performance improvement of the new HR over the existing BR
        when calculating one step $\ham \ket{\psi_n}$ of the power iteration
        excluding normalization.
        HR uses four sites (including impurity) in the bit component.
        Here, \numlist{101; 301} bath sites $n_b$ were tested.
        Left: single excitation.
        Right: single plus double excitation.
    }
    \label{fig:performance-improvement-iteration}
\end{figure}

We can take the measured value at the highest iteration which corresponds to
the calculation of $\ham\ket{\psi_{99}}$ and compare it against the number of
bath sites $n_b$,
as shown in~\cref{fig:ciwf-scaling}.
HR-1 shows a good agreement with a linear fit,
while for HR-2 we can see a scaling which is a little worse than $O(n_b^2)$.

\begin{figure}[htb!]
    \centering
    %! TeX root = ../thesis.tex
\begin{tikzpicture}
    \begin{groupplot}[
            group style={
                    group size=2 by 1,
                    horizontal sep=15mm,
                    ylabels at=edge left,
                },
            width=0.45\textwidth,
            xlabel={$n_b$ (bath sites)},
            ylabel={$t$ (\unit{\second})},
            xtick={100,200,...,300},
            minor xtick={50,150,...,250},
            xminorticks=true,
            legend style={
                    at={(0.05, 0.95)},
                    anchor=north west,
                },
            legend cell align={left},
        ]

        \nextgroupplot[
            xmode=log,
            ymode=log,
            xticklabel=\pgfmathparse{exp(\tick)}\pgfmathprintnumber{\pgfmathresult},
            xticklabel style={
                    /pgf/number format/fixed,
                    /pgf/number format/precision=0,
                },
        ]
        \addplot+[only marks]
        table [x=bath_sites, y=double] {./data/wave_function_time_scaling.csv};
        \addlegendentry{HR-2}
        \addplot+[only marks]
        table [x=bath_sites, y=single] {./data/wave_function_time_scaling.csv};
        \addlegendentry{HR-1}
        \addplot
        [
            domain=50:300,
            samples=10,
            dashed,
        ]
        {3e-6 * x^2}
        coordinate [pos=0.9] (quadratic);
        \node[inner sep=3mm, below] at (quadratic.south) {$O(n_b^2)$};

        \nextgroupplot
        \pgfplotsset{cycle list shift=1}
        \addplot+[
            only marks,
            forget plot,
        ]
        table [x=bath_sites, y=single] {./data/wave_function_time_scaling.csv};
        \addplot
        [
            domain=50:300,
            samples=10,
            dashed,
        ]
        {9.3e-7 * x + 5e-5}
        coordinate [pos=0.9] (linear);
        \node[inner sep=5mm, below] at (linear.south) {$O(n_b)$};
    \end{groupplot}
\end{tikzpicture}

    \caption{
        Time to calculate $\ham\ket{\psi_{99}}$ in the
        hybrid representation for singles (HR-1)
        and singles plus doubles (HR-2).
    }
    \label{fig:ciwf-scaling}
\end{figure}
